\section{Xây dựng mô hình YOLO baseline}

\subsection{Dataset}

Tập dữ liệu sử dụng là \cbs{tên và nguồn dataset}, gồm \cbs{số ảnh} ảnh và
\cbs{số lớp} lớp đối tượng. Dữ liệu được chia thành train, validation và test theo
tỷ lệ \cbs{tỷ lệ chia}. Nếu nhiều frame đến từ cùng video hoặc cùng hiện trường,
các frame liên quan được giữ trong cùng một tập để tránh rò rỉ ngữ cảnh.

Trước huấn luyện, nhãn được kiểm tra về class id, bounding box vượt biên, ảnh lỗi,
đối tượng quá nhỏ và mất cân bằng lớp. Báo cáo cần bổ sung thống kê số instance
theo lớp, phân bố kích thước hộp và ví dụ dữ liệu đại diện.

\subsection{Tiền xử lý dữ liệu}

Ảnh được chuyển về kích thước \cbs{image size}, áp dụng letterbox và chuẩn hóa
theo cấu hình của implementation YOLO. Các phép augmentation gồm
\cbs{liệt kê augmentation thực tế}. Cấu hình tiền xử lý được lưu để dùng lại khi
export và trong Edge Runtime; mọi khác biệt về thứ tự kênh hoặc scale có thể làm
kết quả suy luận sai lệch.

\subsection{Huấn luyện}

Baseline được huấn luyện với optimizer \cbs{optimizer}, learning rate
\cbs{giá trị}, batch size \cbs{giá trị} trong \cbs{số epoch} epoch. Seed, phiên
bản CUDA, framework và commit mã nguồn được lưu cùng checkpoint. Checkpoint
tốt nhất được chọn theo validation $mAP@50{:}95$, không chọn theo tập test.

\subsection{Đánh giá baseline}

\begin{table}[H]
\centering
\caption{Kết quả chất lượng của mô hình baseline}
\label{tab:baseline-quality}
\resizebox{\textwidth}{!}{%
\begin{tabular}{lccccc}
\toprule
\textbf{Model} & \textbf{Precision} & \textbf{Recall} & \textbf{mAP@50} &
\textbf{mAP@50:95} & \textbf{Size (MB)}\\
\midrule
Baseline FP32 & \cbs{} & \cbs{} & \cbs{} & \cbs{} & \cbs{}\\
\bottomrule
\end{tabular}}
\end{table}

Ngoài chỉ số tổng hợp, cần phân tích AP theo lớp và confusion matrix để nhận diện
những lớp yếu. Kết quả này là đường chuẩn để xác định phần accuracy bị mất do
pruning, quantization hoặc chuyển đổi runtime.